\section{病、饥与秩序：危机中的救济和公共责任}

\subsection{灾情如何进入文书}

南宋官府知道一场饥荒、洪水或疫病，往往先通过地方报告、粮价、逃户、死亡与道路损坏的消息。报告把复杂的灾情变成可上呈的类别：某州请求蠲免，某仓请发粮，某段堤防须修，某位官员因延误受责。这种分类是救济能够启动的条件，也会筛掉难以证明、难以登记或尚未被官府注意到的痛苦。

诏令中的赈济、蠲免和禁约不应被轻视。它们显示国家承认维持生计与维持秩序相连，也说明官府试图调动仓储、运输和地方人员。可是，颁令与到达之间有很长的路。粮食要从有余的地方运来，受灾者要被登记，地方官要在军需和救济之间分配有限资源；河道、道路、价格和胥吏的能力都会改变结果。救济数字因此首先是一次行政行动的记录，而不是灾民总数的透明镜子。

\subsection{医疗、寺院与地方互助}

疾病在史料中常不如战役和政令清晰。医书、方书、碑刻、笔记和寺院材料可以显示治疗知识、施药、祈祷或安葬的某些做法，却很少让我们精确测量一次疫病如何穿过整个地区。医疗不是单一职业或单一机构的事务：家族照料、药材市场、僧道活动、地方官的施治和有技术的医者都可能参与。

寺院与社区的施食、施药、收葬有时成为灾时可见的互助。功德记所强调的通常是施主的善行和寺院的名望，它不能证明所有饥病者得到相同帮助。反过来，材料的沉默也不能被理解为社区毫无互助。较谨慎的写法是把寺院看作多种救济渠道之一：它有资源、有门槛，也受地方政治和经济条件限制。

\subsection{价格、治安与民生}

粮价上升会把灾害、战争和财政压力带进每天的购买。官府可能禁止囤积、平粜或调运，商人则根据货源、风险和信用作出判断，家庭会改变消费、借贷或迁居计划。一项价格记录能提示紧张发生，不能单独说明原因或每户承受程度。价格既受收成和运输影响，也会受军需、货币、市场预期和地方权力影响。

治安同样不是与民生分开的专题。饥饿、失业、逃亡和军队往来可能增加冲突，地方官的巡逻、拘捕和禁令是对秩序压力的回应。官书中的盗贼与骚动标签，常把复杂生计压缩为威胁；读者必须问，在被称为“乱”的行动前，是否已有征发、欠粮、道路中断或家庭离散。不能因此把所有治安案件都解释为社会反抗，但也不能只接受官府的分类。

\subsection{危机后的公共世界}

修堤、重开仓场、补葺寺院和恢复市集，是危机后公共世界被重新组织的时刻。谁出钱、谁出工、谁获得表彰，都会影响地方权威。官员、家族、商人和寺院可能在同一工程中合作，也可能围绕费用和功劳竞争。碑刻与文集使其中较有名的参与者可见，普通劳役者往往仍无名；这种差别正是公共设施如何建立的不平等部分。

南宋的长期韧性并不在于从不遭遇病、饥或灾害，而在于官府和社会能在许多地方反复组织救济、运输、医疗和修复。它的局限也在于这些能力分布不均，并会被战争和财政需求反复耗损。到宋末多条水陆线路相继受压时，救济与秩序维护愈发难以同军需分开；国家和地方所能动员的资源，必须在守城、养民和维持市场之间作出代价高昂的选择。
